\section{Metodología propuesta}

\subsection{Diseño general}

\begin{figure}[H]
\centering
\resizebox{\textwidth}{!}{%
\begin{tikzpicture}[
    node distance=0.8cm and 0.9cm,
    box/.style={draw=uvblue, rounded corners, fill=softgray, align=center, minimum height=1.0cm, text width=3.2cm},
    arrow/.style={-{Latex[length=2.5mm]}, thick, draw=darkgray}
]
\node[box] (data) {Estímulo y señal BOLD\\bloque de entrenamiento};
\node[box, right=of data] (pinn) {PINN inversa\\estados $s,f,v,q$\\parámetros $\epsilon,\tau,\alpha$};
\node[box, right=of pinn] (ode) {Integración ODE independiente\\con parámetros estimados};
\node[box, right=of ode] (test) {Predicción del bloque retenido\\métricas fuera de muestra};
\draw[arrow] (data) -- (pinn);
\draw[arrow] (pinn) -- (ode);
\draw[arrow] (ode) -- (test);
\end{tikzpicture}%
}
\caption{Flujo de la PINN inversa con evaluación fisiológica fuera de muestra.}
\label{fig:method_flow}
\end{figure}

La formulación final se diseñó después de un piloto en el que la red aproximaba simultáneamente toda la trayectoria temporal. Aunque los residuos físicos eran bajos, la red podía representar amplitudes diferentes en bloques repetidos. Para evitar esta solución blanda, la formulación definitiva separó la inferencia de parámetros de la predicción: la PINN se ajusta en una ventana de entrenamiento y la evaluación se realiza mediante integración numérica independiente de las ODE.

\subsection{Función de pérdida}

La pérdida total se definió como

\begin{equation}
\mathcal{L}=\lambda_d\mathcal{L}_{\mathrm{datos}}+
\lambda_p\mathcal{L}_{\mathrm{física}}+
\lambda_i\mathcal{L}_{\mathrm{inicial}}+
\lambda_t\mathcal{L}_{\mathrm{terminal}}.
\end{equation}

El término de datos penaliza la diferencia entre BOLD observado y predicho. La pérdida física suma los residuos cuadráticos de las cuatro ecuaciones diferenciales. Las condiciones iniciales imponen homeostasis: $s=0$ y $f=v=q=1$. El término terminal favorece el retorno hacia la línea base al final de la ventana.

\subsection{Parametrización acotada}

Los parámetros se representaron mediante transformaciones sigmoides para garantizar los límites:

\begin{align}
0{,}03 &\leq \epsilon \leq 0{,}35, \\
0{,}50 &\leq \tau \leq 2{,}00, \\
0{,}15 &\leq \alpha \leq 0{,}60.
\end{align}

Estas cotas estabilizan la optimización y evitan regiones no fisiológicas, pero también condicionan la interpretación de los parámetros resultantes.
